%! Characteristics of Linear Functions — Unit 1, Lesson 1.0 — Group Activity
% Gradual release: the "you do together" block, 25 minutes, groups of 3–4.
% Untiered (user direction, 2026-08-31) — every group works the same task, with
% Part 4 as the modelling stretch. The vocabulary was built in the Guided Notes.
\documentclass[10pt]{article}
\usepackage{algebra2-article}
\usepackage{algebra2-boxes}
% Temperature graph: hours h on the horizontal axis (0..7), temperature T on the vertical (-8..8).
% The h-axis is drawn at T=0 so the crossing is the point students circle.
\newcommand{\tempgrid}{%
  \draw[step=1, linegray!40, very thin] (0,-8) grid (7,8);
  \draw[->, semithick, charcoal] (0,0)--(7.6,0) node[right, font=\scriptsize]{$h$};
  \draw[->, semithick, charcoal] (0,-8.6)--(0,8.8) node[above, font=\scriptsize]{$T$ ($^\circ$C)};
  \foreach \x in {1,...,7} \node[charcoal, font=\tiny, below] at (\x,0){\x};
  \foreach \y in {-8,-6,-4,-2,2,4,6,8} \node[charcoal, font=\tiny, left] at (0,\y){\y};}

\begin{document}

\pageheader{Unit 1, Lesson 1.0}{Group Activity}

\vspace{0.08in}

\begin{headlinebox}{forestbg}
\small\textbf{Freezing Point.} On a cold weekend the number that matters most is \emph{zero}: above it
the roads are wet, below it they are ice. Work with your group. Use the names from the notes ---
\textbf{slope}, \textbf{$y$-intercept}, \textbf{zero}, \textbf{increasing/decreasing},
\textbf{positive/negative} --- and for every answer be ready to show \emph{where you see it on the
graph}.
\end{headlinebox}

\vspace{0.06in}

\begin{scenariobox}[1. Saturday Morning]{forest}
At 6:00 a.m.\ on Saturday the temperature in Blacksburg was $-6\,^\circ$C. It rose steadily,
$2\,^\circ$C every hour, all morning (until noon). Let $h$ be the number of hours after 6:00 a.m.
and $T$ the temperature in $^\circ$C.

\begin{enumerate}[label=\textbf{\alph*.}, leftmargin=1.8em, itemsep=5pt]
  \item Complete the table.
    \begin{center}
    \renewcommand{\arraystretch}{1.4}
    \begin{tabular}{|c|c|c|c|c|c|c|c|}
    \hline
    $h$ & $0$ & $1$ & $2$ & $3$ & $4$ & $5$ & $6$ \\ \hline
    $T$ & $-6$ & $-4$ & \blank{0.9cm} & \blank{0.9cm} & \blank{0.9cm} & \blank{0.9cm} & \blank{0.9cm} \\ \hline
    \end{tabular}
    \end{center}
  \item Write a rule for the temperature $h$ hours after 6:00 a.m. \quad
        $T(h) = $ \blank{3.4cm}
\end{enumerate}

\begin{minipage}[t]{0.58\linewidth}
\vspace{0pt}
\begin{enumerate}[label=\textbf{\alph*.}, start=3, leftmargin=1.8em, itemsep=5pt]
  \item Circle the \textbf{$y$-intercept} on the graph: $(\blank{0.7cm},\ \blank{0.7cm})$.
        In the story it tells you:
        \writelines{2}
  \item Circle the \textbf{zero}: $(\blank{0.7cm},\ \blank{0.7cm})$. What time of day is that,
        and why is that moment special?
        \writelines{2}
\end{enumerate}
\end{minipage}\hfill
\begin{minipage}[t]{0.38\linewidth}
\vspace{0pt}\centering
\begin{tikzpicture}[x=0.55cm, y=0.24cm]
  \tempgrid
  \draw[very thick, forest] (0,-6)--(7,8);
  \fill[charcoal] (0,-6) circle (2.8pt);
  \fill[charcoal] (3,0) circle (2.8pt);
\end{tikzpicture}
\end{minipage}

\begin{enumerate}[label=\textbf{\alph*.}, start=5, leftmargin=1.8em, itemsep=5pt]
  \item The \textbf{slope} is \blank{1.2cm} $^\circ$C per hour. Circle where that number sits in
        your rule, then say how you \emph{see} it on the graph from one hour to the next.
        \writelines{2}
  \item Is $T$ \textbf{increasing} or \textbf{decreasing}? \blank{2.4cm} \;
        How does the \emph{graph} show it (not the story)? \blank{4.0cm} \\[3pt]
        $T$ is \textbf{negative} when \blank{2.4cm} \qquad
        $T$ is \textbf{positive} when \blank{2.4cm}
\end{enumerate}
\end{scenariobox}

\boxguard[22]
\begin{scenariobox}[2. Sunday Evening]{navy}
At 4:00 p.m.\ on Sunday it was $6\,^\circ$C, and the temperature \emph{fell} steadily, $2\,^\circ$C
every hour, through the evening. Now $h$ counts the hours after 4:00 p.m.

\begin{minipage}[t]{0.56\linewidth}
\vspace{0pt}
\begin{enumerate}[label=\textbf{\alph*.}, leftmargin=1.8em, itemsep=5pt]
  \item Write a rule: \; $T(h) = $ \blank{3.0cm} \\[2pt]
        Slope: \blank{1.4cm} \quad Increasing or decreasing? \blank{2.4cm}
  \item The $y$-intercept is $(\blank{0.7cm},\ \blank{0.7cm})$ and the zero is
        $(\blank{0.7cm},\ \blank{0.7cm})$. What does each one tell you about Sunday evening?
        \writelines{2}
\end{enumerate}
\end{minipage}\hfill
\begin{minipage}[t]{0.40\linewidth}
\vspace{0pt}\centering
\begin{tikzpicture}[x=0.55cm, y=0.24cm]
  \tempgrid
  \draw[very thick, navy] (0,6)--(7,-8);
  \fill[charcoal] (0,6) circle (2.8pt);
  \fill[charcoal] (3,0) circle (2.8pt);
\end{tikzpicture}
\end{minipage}

\begin{enumerate}[label=\textbf{\alph*.}, start=3, leftmargin=1.8em, itemsep=5pt]
  \item At 5:00 p.m.\ ($h=1$) the temperature is \blank{1.1cm} $^\circ$C. Is that number
        \textbf{positive} or \textbf{negative}? \blank{2.0cm} \; At that moment, is the
        temperature \textbf{rising} or \textbf{falling}? \blank{2.0cm} \\[3pt]
        So: can a temperature be \emph{positive} and \emph{falling} at the same time? Explain what
        each word is describing.
        \writelines{2}
  \item $T$ is \textbf{positive} when \blank{2.6cm} \qquad
        $T$ is \textbf{negative} when \blank{2.6cm}
  \item \textbf{The story, then the math.} In the story, which hours actually make sense as inputs?
        \blank{3.8cm} \\[3pt]
        Now \emph{forget the story}: the rule $T(h)=6-2h$ accepts any number $h$, so its
        \textbf{domain} is \blank{2.6cm} and its \textbf{range} is \blank{2.6cm}.
        Why are those two answers different?
        \writelines{2}
\end{enumerate}
\end{scenariobox}

\boxguard[16]
\begin{scenariobox}[3. Side by Side]{forest}
Look at both graphs together. They have the \emph{same} zero at $h=3$.
\begin{enumerate}[label=\textbf{\alph*.}, leftmargin=1.8em, itemsep=5pt]
  \item Saturday is negative \blank{2.2cm} $h=3$ and positive \blank{2.2cm} it.
        Sunday is the other way around.
  \item Their slopes are \blank{1.4cm} and \blank{1.4cm}. What do the two \emph{signs} tell
        you, and what does the fact that both are $2$ in size tell you?
        \writelines{2}
  \item Write one sentence, using the words \emph{zero} and \emph{sign}, that explains why neither
        line can switch from negative to positive anywhere except at $h=3$.
        \writelines{2}
\end{enumerate}
\end{scenariobox}

\boxguard[22]
\begin{scenariobox}[4. Model It --- The Car-Wash Card]{navy}
A car-wash gift card starts with \$40, and each wash costs \$8, so the balance after $w$ washes is
$B(w) = 40 - 8w$.
\begin{enumerate}[label=\textbf{\alph*.}, leftmargin=1.8em, itemsep=5pt]
  \item What does the slope $-8$ tell you about the card? And the $y$-intercept $(0,40)$?
        \writelines{2}
  \item Find the zero of $B$ by solving $B(w)=0$.
    \begin{work}
      40 - 8w &= 0 \\
          -8w &= -40 \\
            w &= 5
    \end{work}
    In the story, that zero means:
    \writelines{2}
  \item Which inputs $w$ make sense here? \blank{3.4cm} \; Is $B$ increasing or decreasing
        on them? \blank{2.4cm}
  \item Suppose each wash cost \$10 instead of \$8. Would the card run out \emph{sooner} or
        \emph{later}? \blank{2.0cm} \; On the graph, what changes --- where the line
        \emph{starts}, how \emph{steep} it is, or both? Say why.
        \writelines{2}
\end{enumerate}
\end{scenariobox}

\end{document}
